% Draft: case study of an assembly attribution error in a released benchmark
% (GraphPro "E. coli" = E. coli O104:H4 pAA plasmid TSSs).
% Target: Chapter 4 (data provenance) or Chapter 9 (Discussion, benchmark
%   reliability).
% Status: draft v1, 2026-11-02. References checked against Crossref DOIs.
% Project sources: docs/graphpro_derived_artifact_impact.md,
%   scripts/audit_graphpro_ecoli_paa_source.py,
%   configs/graphpro_ecoli_paa_source_audit_v1.yaml, F-013, F-115, F-117.

\section{Case Study: An Assembly Attribution Error in a Released Benchmark}\label{draft:paa-case}

\draftmeta{Chapter~4 (data provenance) or Chapter~9 (benchmark reliability)}
{Draft v1 --- facts audited; exclusion decided by the project owner (log R-040)}
{\texttt{docs/graphpro\_derived\_artifact\_impact.md};
\texttt{configs/graphpro\_ecoli\_paa\_source\_audit\_v1.yaml}; log R-040; F-013, F-115, F-117}

\subsection{What was found}

The GraphPro benchmark \citep{zhang2024graphpro} releases an
\textit{Escherichia coli} promoter set without a construction script. Its
released class files already showed a label problem: 1,345 sequences occur
identically in both the positive and the negative file (F-013). The project
therefore attempted to rebuild the set from independent sources, starting
from the Promotech training inventory \citep{chevezguardado2021promotech}, whose
\texttt{ECOLI.bed} file declares the chromosome of \textit{E.~coli} K-12
MG1655 (\texttt{NC\_000913.3}).

A row-level audit traced the 278 rows of that file to Supplementary Dataset~S1
of \citet{berger2016paa}, a dRNA-seq study of the primary transcriptome of the
74-kb pAA plasmid of \textit{E.~coli} O104:H4. All 278 rows match the published
table in order, position, strand and locus, and all 278 published 51-nt windows
match the plasmid reference \texttt{NC\_018666.1}. None of them comes from the
chromosome named in the file. The coordinates are correct; the replicon they
are attributed to is not.

\subsection{Consequences}

Windows extracted at these coordinates from \texttt{NC\_000913.3} are unrelated
chromosomal sequence, so the nominal \textit{E.~coli} reconstruction, its paired
controls, its split, and every distribution, similarity and dependency result
derived from it were invalidated. They are preserved only as provenance. The
error cannot be repaired by renaming the extracted sequences: a corrected
dataset must be extracted again from \texttt{NC\_018666.1} under a new identity,
with every downstream audit repeated. The source files were not edited.

Even with the correct replicon, the source does not define a conventional
promoter set. Under the classification of \citet{sharma2010primary}, the
248 distinct TSSs comprise 51 gene-associated (primary or secondary), 58
internal, 111 antisense and 47 orphan sites, and 100 sites are antisense only
(F-115). Applying the gene-associated rule of Section~\ref{draft:tss-class-selection}
would leave 51 positives.

\subsection{Lessons for benchmark construction}

\paragraph{Coordinates do not establish identity.}
A file's declared accession is metadata, not evidence. Here, sequence-level
agreement with the primary publication was required to detect that the
declared chromosome was wrong. Verifying every source window against its
assembly before extraction is cheap compared with the cost of an invalid
dataset propagating into training, partitioning and dependency analysis.

\paragraph{Errors propagate silently through derived resources.}
The released benchmark, the intermediate inventory and the reconstruction each
looked internally consistent. Only an audit that returned to the primary table
exposed the problem. This mirrors the broader observation that flaws in data
preparation, such as leakage, are a widespread source of overoptimistic results
in machine-learning-based science \citep{kapoor2023leakage}.

\paragraph{The name of a dataset is a claim.}
Labelling a set of plasmid TSSs, most of them non-canonical, as an
``\textit{E.~coli} promoter'' benchmark changes what a model trained on it is
said to recognise. Reconstruction therefore keeps the source designation
verbatim, records the verified organism and replicon separately, and gives a
corrected resource a new identity rather than inheriting the old name.

\subsection{Status}

The GraphPro \textit{E.~coli} set was excluded from the reconstructed line, and
no pAA dataset was built, exploratory or confirmatory (log R-040). Under the
gene-associated rule only 51 positives would remain, too few for a test panel.
Space for genomic controls was not the reason: even with the full exclusion
radius the plasmid leaves 453 non-overlapping slots per strand against 124
needed (F-117). The released files, the nominal reconstruction and the source
audit are kept as provenance; nothing was deleted or relabelled.
